\section{文化与思想：礼不是装饰，是政治语言}

西周留下的“礼乐文明”常被后人说得很抽象。放回当时，它首先是一套让不同身份的人知道彼此位置的公共语言：谁能祭祀，谁能接受册命，谁在朝聘时站在哪里，谁的功劳可以铸在器物上。这些程序既表达敬意，也划出权利和义务的边界。

金文之所以重要，是因为它让我们看见这种语言的实际使用：一次赏赐、一次征伐、一段家族记忆被铸进器物，便成为后代仍可调用的政治凭据。后世儒家把周礼理解为理想秩序，是一种重要的历史记忆；但西周人实际面对的，是如何用礼把权力、资源和宗族关系暂时安顿下来。

\begin{textwindow}[后来的孔子怎样回望周]
《论语》有“周监于二代，郁郁乎文哉，吾从周”。这句话适合放在西周章中，不是因为孔子生活在西周，而恰恰因为他生活在西周已亡数百年后的春秋末期。它让我们看见：周既是一个已经消失的政治秩序，也是后来人用来批评现实、想象秩序的一面镜子。
\end{textwindow}
